% ======================================================
%  PART I — VẤN ĐỀ CẦN GIẢI QUYẾT
% ======================================================
\section{Các vấn đề cần giải quyết}

Dựa trên kết quả phỏng vấn và quan sát người dùng, nhóm xác định ba vấn đề nghiêm trọng cần được giải quyết như đã liệt kê như bên dưới. Với mỗi vấn đề, nhóm mô tả chúng dưới góc nhìn của người dùng, đưa ra giải thích phù hợp nhất về vấn đề xảy ra, và đề xuất một số phương án giải quyết.

% ======================================================
%  PROBLEM 1: TRÌ HOÃN (LAG) KHI XEM TRỰC TIẾP
% ======================================================
\subsection{Vấn đề 1: Trì hoãn (\textit{lag}) khi xem trực tiếp}

Trong quá trình phỏng vấn, người được phỏng vấn xác định rằng trì hoãn (\textit{lag}) là ``vấn đề gây khó chịu lớn nhất'' khi xem nội dung thể thao trực tuyến. Người này cho biết ``xem thể thao mà bị lag sẽ khiến họ muốn tắt ngay lập tức''. Đây là phản hồi mạnh mẽ nhất trong tất cả các vấn đề được đề cập. Người này cũng dõng dạc khẳng định rằng ``Bây giờ có hỏi bao nhiêu người trong lớp cũng thế thôi, ai cũng sẽ trả lời y chang như vậy''.

\begin{itemize}[leftmargin=*, nosep]
    \item \textbf{Video dừng đúng lúc cảm xúc đang dâng cao — khoảnh khắc tệ nhất có thể}: Khi xem trực tiếp, người xem hưng phấn nhất đúng vào những pha quyết định — tình huống sắp ghi bàn, pha bóng căng thẳng sắp phân định kết quả. Cũng chính lúc này, nền tảng lại hay bị trễ nhất do lượng người xem tăng đột biến. Video dừng lại đúng khoảnh khắc cảm xúc đang chuẩn bị lên đến đỉnh: sự hồi hộp không được giải tỏa mà bị bỏ lửng, cảm xúc tụt xuống ngay khi nó lẽ ra phải vỡ òa. Vài giây chờ đợi tưởng chừng nhỏ, nhưng nó cắt đứt đúng cao trào — trải nghiệm mượt mà của một pha bóng bị phá vỡ chỉ vì vài giây trễ.
    \item \textbf{Mất niềm tin vào nền tảng}: Khi trải nghiệm không mượt mà, người dùng mất niềm tin vào chất lượng tổng thể của ứng dụng và có xu hướng chuyển sang nền tảng khác.
\end{itemize}

Nền tảng phát trực tuyến thể thao cần ưu tiên tối ưu hóa hiệu suất để giảm thiểu độ trễ, đặc biệt trong các khung giờ cao điểm khi lượng người xem tăng đột biến. Trong hầu hết trường hợp, đây là yếu tố then chốt quyết định người dùng có ở lại nền tảng hay không. Ngoại lệ duy nhất là khi nội dung có tính độc quyền cao đến mức người dùng không có lựa chọn khác.

% ======================================================
%  PROBLEM 2: GIAO DIỆN QUÁ NHIỀU HÌNH ẢNH, THIẾU VĂN BẢN
% ======================================================
% \subsection{Vấn đề 2: Giao diện quá nhiều hình ảnh, thiếu thông tin văn bản}
%
% Người được phỏng vấn nhận xét rằng giao diện ``hiển thị quá nhiều hình ảnh trong khi thông tin bằng văn bản còn hạn chế, làm người dùng khó nhận biết nhanh nội dung cần xem''. Trên màn hình đầu tiên, người dùng phải đối mặt với hàng loạt thumbnail, banner và carousel mà gần như không có chú thích hay mô tả nào giúp phân biệt nội dung.
%
% \begin{itemize}[leftmargin=*, nosep]
%     \item \textbf{Văn bản thể hiện được những thông tin mà hình ảnh không thể:} Trong trường hợp người dùng ít khi theo dõi thể thao, khi vào trang chủ VTVPrime và nhìn thấy một thumbnail của cầu thủ, họ sẽ không biết đó là cầu thủ nào, thuộc đội nào, đang đấu giải nào — ảnh chỉ gợi cảm xúc chứ không truyền tải thông tin cụ thể. Một dòng văn bản đơn giản như ``Man United vs Liverpool — LIVE — 2--1'' cung cấp đủ thông tin ngay lập tức. Hình ảnh rất tốt để thu hút ánh nhìn, nhưng phải có văn bản bổ trợ thì người dùng mới ``hiểu nhanh'' được nội dung.
%     \item \textbf{Thiếu văn bản buộc người dùng phải cố gắng nhớ ra nội dung là gì thay vì nhận ra ngay lập tức:} Nguyên tắc ``tối ưu khả năng nhận diện, giảm khả năng gợi nhớ'' (Recognition over Recall) phân biệt hai cách não xử lý thông tin: nhận diện (nhìn thấy, nhận ra ngay) và gợi nhớ (phải lục lại trí nhớ). Khi ảnh thu nhỏ không có văn bản đi kèm, người dùng rơi vào trạng thái gợi nhớ. Trong tâm trí họ xuất hiện những câu hỏi như ``trận này là trận nào nhỉ?''. Ngược lại, khi có văn bản bổ sung thông tin, não chuyển sang chế độ nhận diện — thấy là hiểu ngay, không cần nỗ lực để nhớ ra một thông tin nào đó nữa. Giao diện tốt nên giúp người dùng luôn ở chế độ nhận diện, không phải gợi nhớ.
% \end{itemize}
%
% Giao diện cần cân bằng giữa hình ảnh và văn bản: giữ ảnh thu nhỏ để thu hút sự chú ý của người dùng, nhưng bổ sung thông tin dạng văn bản rõ ràng (tên trận đấu, tỷ số,\dots) để người dùng có thể nhận biết nhanh nội dung cần xem. Đây là vấn đề đặc biệt quan trọng trên điện thoại khi diện tích màn hình hạn chế.

% ======================================================
%  PROBLEM 2: QUÁ TẢI THÔNG TIN CHO NGƯỜI DÙNG MỚI
% ======================================================
\subsection{Vấn đề 2: Quá tải thông tin đối với người dùng mới}

Người dùng mới cảm thấy giao diện ``khá rối'' vì ``có quá nhiều thông tin xuất hiện ngay trên màn hình đầu tiên''. Ngay lần đầu mở ứng dụng, họ phải tiếp nhận cùng lúc hàng loạt banner, carousel và thumbnail khác nhau, dẫn đến không biết nên bắt đầu từ đâu.

\begin{itemize}[leftmargin=*, nosep]
    \item \textbf{Quá tải bộ nhớ làm việc (\textit{working memory})}: Bộ nhớ làm việc chỉ chứa được khoảng $7 \pm 2$ khối thông tin (chunk) cùng lúc, với thời lượng lưu giữ khoảng 10 giây~\cite{miller1956}. Khi giao diện hiển thị quá nhiều thành phần, người dùng mới bị quá tải nhận thức — bộ nhớ làm việc không đủ dung lượng để xử lý tất cả.
    \item \textbf{Sự chú ý của người dùng bị phân tán}: Khi mở ứng dụng, người dùng phải lia mắt đến hàng loạt banner, carousel và thumbnail trên giao diện. Mỗi loại thành phần này đều khá tương đồng với nhau về hình dáng và dường như không có sự phân cấp rõ ràng, khiến não bộ khó tập trung vào một nội dung cụ thể nào đó.
\end{itemize}

Cách tiếp cận phù hợp là chỉ hiển thị những nội dung quan trọng nhất và thêm cơ chế để người dùng xem thêm nội dung khi họ cần. Giảm bớt số lượng thông tin hiển thị ban đầu giúp người dùng mới không bị choáng ngợp, đồng thời vẫn giữ được khả năng tiếp cận đầy đủ nội dung cho người dùng có kinh nghiệm.